\documentclass[10pt]{article}
\usepackage[margin=0.8in]{geometry}
\usepackage{amsmath,booktabs}
\setlength{\parindent}{0pt}
\setlength{\parskip}{6pt}
\begin{document}
\begin{center}
{\Large\bfseries Question 6: Batch Training versus SGD}\par
Observed behavior and explanation
\end{center}

\section*{Training conditions}
Both models start from weights $[1,-1]$ and train for 300 epochs with learning
rate $0.01$ and momentum $0.5$. Batch training updates once per epoch; SGD
updates after each point, with shuffled order using seed 0. This gives 300
batch updates and 900 SGD updates. Both comparison curves show full-data MSE
after each epoch.

\section*{Observed behavior}
With the same initialization and optimizer settings, batch training oscillates
initially, then settles at a nearly constant loss. Per-point SGD has a
much noisier curve, including large fluctuations late in training.
\begin{center}
\begin{tabular}{lrr}
\toprule
Method & Parameter updates & Final full-data MSE\\
\midrule
Batch gradient descent & 300 & 5.476191\\
Per-point SGD & 900 & 5.776175\\
\bottomrule
\end{tabular}
\end{center}
\section*{Why this happens}
Batch training averages the gradients from all three data points before each
update. This balances their competing effects, producing a more stable loss
curve after the initial oscillations.

SGD updates using one point at a time. An update that improves one point can
worsen the others, so the full-data loss fluctuates. Shuffling the points changes
the order of these updates. The constant learning rate keeps the updates from
shrinking automatically, and momentum carries previous updates forward, which
can amplify overshooting.
\end{document}
